%% Lederman Fellowship -- Outreach Statement
%% Sean Patrick MacBride
%%
%% Compile with: pdflatex -> pdflatex (for \pageref{LastPage})
%% Style deliberately matches researchStatement.tex / coverLetter/letter.tex
%% so that the three application documents read as one packet.

\documentclass[11pt]{article}
\renewcommand{\baselinestretch}{1}
\usepackage[
  a4paper,
  margin=0.85in,
  headsep=3pt,
]{geometry}
\usepackage[most]{tcolorbox}
\usepackage{xcolor}
\usepackage{fancyhdr}
\usepackage{lastpage}
\usepackage[inline]{enumitem}
\setlist[itemize]{noitemsep, topsep=0pt, leftmargin=1.2em, itemsep=-0.1cm}
\usepackage{marginnote}
\usepackage[
  breaklinks,
  colorlinks=true,
  linkcolor=blue,
  urlcolor=blue,
  citecolor=blue,
  pdftitle={Sean Patrick MacBride - Outreach Statement},
  pdfauthor={Sean Patrick MacBride},
  unicode
]{hyperref}

\newcommand{\TODO}[1]{%
  \marginnote{\color{red}\textbf{TODO}}%
  {\color{red}#1}%
}

%% ---------------------------------------------------------------
%% Personal-info macros (kept in the same style as researchStatement.tex
%% so a single find/replace keeps every document in sync)
%% ---------------------------------------------------------------
\newcommand{\yourname}{Sean Patrick MacBride}
\newcommand{\youremail}{\href{mailto:sean.macbride@physik.uzh.ch}{sean.macbride@physik.uzh.ch}}
\newcommand{\yourweb}{\href{https://seanmacb.github.io}{seanmacb.github.io}}
\newcommand{\soptitle}{Outreach Statement}

\pagestyle{fancy}
\fancyhf{}
\fancyhead[C]{%
  \footnotesize\rmfamily
  Web: \textcolor{blue}{\yourweb}\quad
  \yourname\ -- Outreach Statement \quad
  E-mail: \textcolor{blue}{\youremail}}
\fancyfoot[C]{Page \thepage\ of \pageref{LastPage}}

%% Section-header macro (matches researchStatement.tex)
\newcommand{\statement}[1]{\par\medskip
  \textcolor{blue}{\textbf{#1:}}%
}

\begin{document}
\hrule
\begin{center}
    \Large\soptitle\ -- \yourname
\end{center}
\vspace{0pt}
\hrule
\vspace{0.5pt}
\hrule height 0.5pt
\vspace{3pt}
\thispagestyle{plain}

\normalsize
\noindent
Leon Lederman spent as much energy convincing the public that physics belongs to everyone as he spent doing the physics itself, and I applied to this fellowship because I want a career built on that same premise. My own path into astrophysics began at a small liberal-arts telescope, not a national laboratory, and every outreach activity I have taken on since -- from public nights at that telescope to gravitational-wave talks at a city-wide museum night -- has been an attempt to hand someone else the same first look at the sky that hooked me. I would bring that same instinct to Fermilab, using the Lederman Science Center, Fermilab's Office of Education and Public Engagement, and the wider Chicagoland science community as the platform, and multi-messenger, gravitational-wave cosmology as the hook.

\begin{tcolorbox}[
  colback=blue!5!white,
  colframe=blue!75!black,
  boxsep=0.4mm, top=0.4mm, bottom=0.4mm, left=2mm, right=2mm,
  toptitle=0.4mm, bottomtitle=0.4mm, before skip=2.5pt, after skip=3.5pt,
  title=Outreach Highlights
]
  \begin{itemize}
    \item Led public ``Ask a Scientist'' sessions on gravitational-wave cosmology for the University of Zurich's Long Night of Museums (2026), reaching walk-in public audiences of all ages.
    \item Guided tours and press briefings on the LSST Camera for NSF/DOE review panels, government delegations, and the public at Vera C. Rubin Observatory (2025).
    \item Led high-school students through observation and analysis of the 2023 annular and 2024 total solar eclipses, with paired lectures on solar physics (Burlington, VT).
    \item Ran public telescope nights at Wheaton College Observatory for four years, and have taught large introductory physics courses at the University of Michigan and the University of Zurich.
    \item Mentored six undergraduate and master's students since 2023, several of whom have gone on to PhD programs in physics.
  \end{itemize}
\end{tcolorbox}

\statement{Past outreach and teaching experience}
My outreach record spans public astronomy, classroom teaching, and mentorship, and each has shaped how I plan to contribute as a Lederman Fellow. As an undergraduate and early graduate student at Wheaton College's observatory (2016--2020), I ran public telescope nights, helping visitors of every background find and understand objects in the night sky -- my first experience translating a technical instrument into a shared, human experience. That instinct carried into the classroom: I have been the lead or sole instructor for large introductory physics courses at the University of Michigan and, currently, the University of Zurich, and I designed a public-facing art installation in the Michigan physics department documenting the struggles and eventual successes of the department's own graduates, aimed at demystifying what a career in physics actually looks like.

Two experiences map especially closely onto what Lederman Fellows are asked to do. First, during the 2023 annular and 2024 total solar eclipses, I led a group of Vermont high-school students through the physics of the eclipses, preparing and delivering lectures on solar observation both virtually and in person -- a self-contained, standards-adjacent outreach unit not unlike the modules Fermilab's QuarkNet teachers bring back to their own classrooms. Second, in 2025 I was a primary guide for visiting parties at Rubin Observatory as we commissioned the LSST Camera, translating the same instrument for audiences ranging from touring media and dignitaries to NSF and DOE review teams to friends and family -- practice in adjusting one explanation of frontier instrumentation to whatever room I am in, which is exactly the skill Saturday Morning Physics and Quantum lecturers rely on. Most recently, in 2026 I led ``Ask a Scientist'' sessions on observational cosmology and gravitational waves at the Physik-Institut's stand for the City of Zurich's Long Night of Museums, an annual event in which the city's science, art, and culture museums open their doors free of charge -- the closest analogue in my record to a Fermilab Family Open House or Ask-a-Scientist booth, and the format I most want to bring across the Atlantic.

Mentoring has been the most sustained form of outreach in my career. Since 2023 I have mentored six undergraduate and master's students -- Elise Fuller and Tim Fanning at the University of Michigan, and Banan Yamani, Yavuz Gencel, Guandi Zhao, and Alexs Gaute at the University of Zurich -- guiding each through a first research project, from formulating a question to presenting the result. Three are now pursuing physics PhDs of their own (at the University of Zurich, the Paul Scherrer Institute, and CU Boulder) and one is now a controls engineer at Argonne National Laboratory. I see this one-on-one work as outreach with a longer time horizon than a single lecture or open house: convincing one more person that a research career is available to them.

\statement{A plan for outreach as a Lederman Fellow}
Fermilab's Lederman Fellows have a strong recent record of pairing frontier research with public science communication -- the current cohort's Dark Matter Days events, science-trivia nights, and QICK tutorials for the quantum program are the model I would want to extend to gravitational-wave and multi-messenger cosmology, a topic the public already meets through LIGO-Virgo-Kagra detection announcements but rarely gets to explore hands-on. Concretely, I would:

\begin{itemize}
  \item \textbf{Lecture in Saturday Morning Physics and Saturday Morning Quantum.} I would propose and lead sessions on gravitational-wave discovery and standard-siren cosmology -- how a single neutron-star merger, seen in both light and gravitational waves, lets us measure the expansion of the universe -- built directly from my own research and from the public-facing version of that story I already give at Zurich's Long Night of Museums.
  \item \textbf{Run Ask-a-Scientist and Family Open House sessions at the Lederman Science Center.} These map onto the format I have already run successfully in Zurich and the visitor-facing work I did at Rubin Observatory, and I would use them to connect Fermilab's neutrino and dark-matter program to the electromagnetic-follow-up and cosmology work I do with Rubin, the Dark Energy Camera, and the Prime Focus Spectrograph.
  \item \textbf{Bring a gravitational-wave module to QuarkNet teacher workshops.} My experience preparing eclipse lectures for Vermont high-school teachers and students translates directly into a classroom-ready unit on multi-messenger astronomy, giving QuarkNet teachers a way to connect particle- and astro-physics data analysis for their students.
  \item \textbf{Extend mentorship into Chicagoland's own pipeline programs.} Building on my record of mentoring students into PhD programs, I would volunteer as a judge and project mentor for the Illinois Junior Academy of Science's regional expositions and the CPS Student STEM Exhibition \& Symposium, both of which actively recruit scientists with physics and astronomy backgrounds from the Fermilab area.
  \item \textbf{Mentor Fermilab's own summer researchers.} I would take on students through Fermilab's TARGET high-school internship program and the SULI/CCI undergraduate internship program, extending the same one-on-one research mentorship that has already guided six students through their first projects and sent three of them on to physics PhDs of their own.
  \item \textbf{Partner with Chicago's public astronomy venues.} The Adler Planetarium and the Griffin Museum of Science and Industry both run public astronomy programming (the Adler alongside the University of Chicago's own ``Astronomy Conversations'' series) that would let me bring gravitational-wave cosmology to audiences well beyond the lab, complementing Rubin Observatory's own public data releases as they begin in earnest during my fellowship.
\end{itemize}

Taken together, these are not activities I would bolt onto a research fellowship -- they are a continuation of what I have already been doing at every stage of my career, now with the platform, the institutional history, and the surrounding Chicagoland science ecosystem that only a Lederman Fellowship can offer. I would be proud to carry that piece of Leon Lederman's legacy forward.

\end{document}
